% ============================================================
% SECTION 6 — FINANCIAL SAFETY MODEL
% ============================================================

\section{Financial Safety Model}

\begin{tcolorbox}[summarybox={\iconShield[1.5] The "At-Most-Once" Guarantee}]
\color{textdark}
The primary technical mandate of RecoverAI V2 is the \textbf{At-Most-Once} financial execution guarantee. A transaction must never be charged twice, and a refund must never be issued twice. 

Because networks drop packets, processes crash mid-execution, and clients retry aggressively, the system assumes that duplicate execution attempts \textit{will} occur. The Financial Safety Model is a combination of design patterns designed to deterministically block these duplicates.
\end{tcolorbox}

\vspace{0.8cm}

% --- THE THREE PILLARS ---
\subsection{The Three Pillars of Financial Safety}

The model is built on three distinct pillars, each defending a specific boundary of the application.

\begin{center}
\begin{tikzpicture}[
    node distance=0.8cm,
    pillar/.style={rectangle, draw=primary, thick, fill=bglight, minimum width=4.5cm, minimum height=1.2cm, rounded corners=4pt, font=\small\bfseries, text=textdark, drop shadow=black!5, align=center},
    arrow/.style={-{Stealth[scale=1.1]}, thick, draw=textmuted}
]

    \node[pillar] (p1) {1. API Idempotency\\(\filepath{IdempotencyRecord})};
    \node[pillar, right=of p1] (p2) {2. State Concurrency\\(\code{with\_for\_update()})};
    \node[pillar, right=of p2] (p3) {3. ExecutionGuard\\(The Final Boundary)};
    
    \node[font=\scriptsize\itshape, text=textmuted, below=0.2cm of p1, align=center] {Protects the Edge\\(Prevents duplicate API calls)};
    \node[font=\scriptsize\itshape, text=textmuted, below=0.2cm of p2, align=center] {Protects the Database\\(Prevents thread races)};
    \node[font=\scriptsize\itshape, text=textmuted, below=0.2cm of p3, align=center] {Protects the Gateway\\(Prevents unsafe state transitions)};

\end{tikzpicture}
\end{center}

\vspace{0.8cm}

% --- PILLAR 1: API IDEMPOTENCY ---
\subsection{Pillar 1: API Idempotency}

Every financial POST request (payment, recovery, refund) requires an \code{Idempotency-Key} header. The API maps this key, combined with the \code{X-API-Key}, into a dedicated PostgreSQL table (\code{idempotency\_records}).

\begin{tcolorbox}[featurebox={Idempotency Lifecycle}]
\begin{enumerate}
    \item \textbf{Initial Request:} System attempts to insert \code{status='PENDING'} into the DB. If successful, processing begins.
    \item \textbf{Concurrent Retry (Race):} If a second request arrives while the first is processing, the DB unique constraint throws an \code{IntegrityError}. The API catches it and returns \code{409 Conflict: Request already in progress}.
    \item \textbf{Completed Retry:} If the first request finished, it updated the record to \code{COMPLETED} and saved the \code{response\_body}. The API intercepts the retry, skips execution entirely, and returns the cached JSON response with a \code{200 OK} and an \code{Idempotent-Replay: true} header.
\end{enumerate}
\end{tcolorbox}

\vspace{0.8cm}

% --- PILLAR 2: OPTIMISTIC CONCURRENCY CONTROL ---
\subsection{Pillar 2: Optimistic \& Pessimistic Concurrency}

Idempotency protects against clients sending the same key twice. But what if two different background workers try to process the \textit{same} database record simultaneously?

\begin{multicols}{2}
\begin{tcolorbox}[infobox={Pessimistic Locking}]
\textbf{Tool:} \code{with\_for\_update()}\\
\textbf{Use Case:} Reconciliation sweepers.
When the background task sweeps for stuck transactions, it uses a row-level lock. If a second worker is running, PostgreSQL forces it to wait or skip the row, ensuring only one worker parses the transaction.
\end{tcolorbox}

\columnbreak

\begin{tcolorbox}[infobox={Optimistic Locking (OCC)}]
\textbf{Tool:} \code{version} column.\\
\textbf{Use Case:} State Machine transitions.
Every attempt has a \code{version} int. When transitioning from \code{PENDING} to \code{AUTHORIZED}, the SQL reads: \code{UPDATE ... WHERE id=X AND version=Y}. If another thread already transitioned it, \code{version} is now \code{Y+1}, 0 rows are affected, and the application aborts.
\end{tcolorbox}
\end{multicols}

\vspace{0.8cm}

% --- PILLAR 3: EXECUTION GUARD ---
\subsection{Pillar 3: ExecutionGuard (The Final Boundary)}

Even with idempotency and database locks, logical errors in the orchestrator could still trigger a duplicate gateway call. \textbf{ExecutionGuard} is a standalone service injected exactly at the gateway boundary. No component in the system is allowed to call the \code{GatewayInterface} without passing through ExecutionGuard.

\begin{center}
\begin{tikzpicture}[
    node distance=1.2cm,
    actor/.style={rectangle, draw=secondary, thick, fill=bglight, minimum width=2.5cm, minimum height=1cm, rounded corners=4pt, font=\small\bfseries},
    guard/.style={rectangle, draw=primary, thick, fill=primary!15, minimum width=3.5cm, minimum height=4cm, rounded corners=4pt, font=\small\bfseries, align=center},
    gateway/.style={rectangle, draw=textdark, dashed, fill=textdark!10, minimum width=2.5cm, minimum height=4cm, rounded corners=4pt, font=\small\bfseries, align=center},
    arrow/.style={-{Stealth[scale=1.1]}, thick, draw=textdark},
    redcross/.style={draw=danger, ultra thick, -, align=center}
]

    % Actors
    \node[actor] (orch) {Orchestrator};
    \node[actor, below=1cm of orch] (refund) {Refund Service};
    
    % Guard
    \node[guard, right=1.5cm of orch, yshift=-1cm] (eg) {ExecutionGuard\\(Safety Checkpoint)};
    
    % Gateway
    \node[gateway, right=1.5cm of eg] (gw) {Payment\\Gateway};
    
    % Safe paths
    \draw[arrow] (orch.east) -- (eg.150);
    \draw[arrow] (refund.east) -- (eg.210);
    \draw[arrow, draw=success, line width=1.5pt] (eg.east) -- node[above, font=\scriptsize\bfseries, text=success] {Valid} (gw.west);
    
    % Internal checks inside EG
    \node[font=\scriptsize\itshape, text=textdark, align=left, yshift=-0.5cm] at (eg.center) {
        $\bullet$ Check prior attempts\\
        $\bullet$ Verify attempt state\\
        $\bullet$ Validate action type\\
        $\bullet$ Inspect parent TX
    };

\end{tikzpicture}
\end{center}

\begin{tcolorbox}[codebox={ExecutionGuard Deep Inspection Logic (\filepath{app/services/execution\_guard.py})}]
\begin{lstlisting}[language=Python]
def execute(self, action: str, attempt: RecoveryAttempt) -> ExecutionResult:
    # 1. Action Allowlist
    if action not in ["RETRY_PAYMENT"]:
        raise ValueError("Invalid action")
        
    # 2. Strict State Enforcement
    if attempt.status != "AUTHORIZED":
        raise ValueError(f"Attempt must be AUTHORIZED, got {attempt.status}")
        
    # 3. Crash-Window Protection (Batch 4.6 fix)
    for prior in transaction.recovery_attempts:
        if prior.id != attempt.id and prior.status in [
            "EXECUTING", "VERIFYING", "UNKNOWN", "SUCCEEDED"
        ]:
            raise ValueError("Prior attempt is active/terminal. Aborting.")
            
    # 4. Safe to execute
    return self.gateway.execute_recovery_action(attempt)
\end{lstlisting}
\end{tcolorbox}

\subsubsection{The Crash-Window Scenario}
Before Batch 4.6, if the system crashed immediately after the gateway returned success but before the database committed the new transaction status, a rebooted worker would see the transaction as \code{NOT\_STARTED} and try again. 

ExecutionGuard's \textbf{Deep Inspection} (Rule 3) prevents this. Even if the parent transaction appears \code{NOT\_STARTED}, the Guard inspects the \textit{child} attempts. Because the gateway call transitioned the child attempt to \code{EXECUTING}, the Guard detects the orphaned active attempt and aborts the new request, perfectly bridging the database-gateway crash window.
